\documentclass[a4paper,11pt]{article}
\usepackage[utf8]{inputenc}
\usepackage[french]{babel}
\usepackage{amsmath}
%\usepackage[osf,sups]{Baskervaldx} % lining figures
%\usepackage[bigdelims,cmintegrals,vvarbb,baskervaldx]{newtxmath} % math font
\usepackage{graphicx}
\usepackage{xcolor}
\usepackage{url}
\usepackage[hidelinks]{hyperref}

% hyperrref configuration: remove ugly boxes and colors, keep urls in blue
\hypersetup{colorlinks,linkcolor=black,citecolor=black,urlcolor={blue!80!black}}


% paragraphs
\setlength{\parindent}{0em}
\setlength{\parskip}{1em}

% margins
%\addtolength{\hoffset}{-3em}
%\addtolength{\textwidth}{6em}
%\addtolength{\voffset}{-3em}
%\addtolength{\textheight}{6em}

% macros
\def\ud{\mathrm{d}}

\begin{document}
\thispagestyle{empty}

{\bf
	Formule de Poisson pour les quasi-cristaux
}

{\bf Encadrants}:\\
Rafael Grompone von Gioi \verb+<rafael.grompone@ens-paris-saclay.fr>+\\
Enric Meinhardt-Llopis \verb+<enric.meinhardt@ens-paris-saclay.fr>+

{\bf Objectif.}
L'objectif final du stage est de construire des figures qui illustrent le la théorie
des quasi-cristaux développée par Yves Meyer le long des 50 dernières années.
Au cours du projet, on rencontrera et on étudiera les concepts proches mais
différents de quasi-cristal, mesure cristalline, ensemble de Meyer, pavage
apériodique.
Un outil essentiel sera une formule de Poisson généralisée
\[
	\sum_{x\in Q}\varphi(x) = \sum_{y\in Q^*}\widehat\varphi(y)
\]
où~$Q$ est un sous-ensemble discret du plan,~$Q^*$ son ``ensemble
dual'' et~$\varphi$ une fonction de Schwartz quelconque.

Avez-vous jamais vu un pavage du plan avec symétrie rotationnelle d'ordre~13~?
Vous trouverez pas de figures sur internet.
Les travaux de Meyer expliquent comment le construire, mais sans jamais le montrer.
Dans ce stage on étudiera le procédé de construction et on produira des belles
images.



{\bf Contexte.}


\medskip

{\footnotesize
\begin{tabular}{ccc}
	\includegraphics[width=0.15\linewidth]{f/georev.png}&
	\includegraphics[width=0.35\linewidth]{f/analogy.jpg}&
	\includegraphics[width=0.3\linewidth]{f/jacomarp.png}\\
	Géodésiques~\cite{rouviere} sur des surfaces &
	Analogie gravité/géométrie &
	Plongements locaux~\cite{darbouxi} d'orbites\\
	&
	(dessin typique mais incorrect!) &
	de Kepler comme géodésiques.
\end{tabular}
}

\medskip



{\bf Déroulement.}
Ce stage s'adresse aux étudiants qui souhaitent approfondir leurs connaissances
en géométrie discrète et analyse harmonique.  Après l'étude du problème de
base, plusieurs ouvertures sont possibles selon les préférences des étudiants:
vers la théorie de nombres et l'approximation diophantienne, vers le traitement
du signal et l'échantillonnage irrégulier, vers l'histoire des pavages
apériodiques, ou même vers la représentation artistique de ceux-ci.


\vspace{-1.5em}
\renewcommand{\refname}{\normalsize Références}
%
\begin{thebibliography}{99}
\vspace{-1em}
{\footnotesize
\bibitem{arnohooke}
	V.I.~Arnold
	{\it Huygens and Barrow, Newton and Hooke}
	Birkhäuser, 1990
}
\end{thebibliography}


\end{document}


% vim:set tw=79 spell spelllang=fr:
